\documentclass[11pt]{article}
\usepackage[T1]{fontenc}
\usepackage{lmodern,amsmath,amssymb,amsthm,booktabs}
\usepackage[margin=1in]{geometry}
\usepackage[colorlinks=true,allcolors=blue]{hyperref}
\newtheorem{theorem}{Theorem}[section]
\newtheorem{lemma}[theorem]{Lemma}
\newtheorem{proposition}[theorem]{Proposition}
\newtheorem{corollary}[theorem]{Corollary}
\theoremstyle{remark}\newtheorem{remark}[theorem]{Remark}
\newcommand{\C}{\mathbb C}
\newcommand{\R}{\mathbb R}
\newcommand{\Q}{\mathbb Q}
\newcommand{\A}{\mathbb A}
\DeclareMathOperator{\trdeg}{trdeg}
\DeclareMathOperator{\Tr}{Tr}
\DeclareMathOperator{\ord}{ord}
\title{Identifying infection stages, infectivity weights,\\and initial states from incidence}
\author{Tatsuhiko Takimoto\\Independent researcher}
\date{Version 1.1.0 --- 2026-11-26}
\begin{document}
\maketitle
\begin{abstract}
A staged SIR model with a common progression rate can assign a different infectivity weight to each infection stage. We prove that its complete noiseless incidence germ in the same physical time and the known total population determine the number of stages, both rates, all infectivity weights, and every initial state, provided both true and competing models have strictly positive rates, weights, and initial states, with the first weight normalized to one. The true and competing stage numbers range over all finite positive integers. In particular, no exceptional positive true points remain in the equal-infectivity model. The proof establishes Zariski density of each positive trajectory after eliminating the removed compartment. A logarithmic first integral forces the algebraic closure to contain the equilibrium line; a boundary argument using relations among Jordan-chain invariants excludes every proper closure of dimension at least two. Polynomial chain consistency then connects the nonlinear incidence observation to a transfer-function identity, which determines the hidden chain. Exact finite checks supplement the proof. For a broader positive compartment feedback class, we also give sharp output-germ jet orders 4, 8, 11, and 15 when one infectious order is one and the other ranges from one to four.
\end{abstract}

\section{The question and the observation}
Pant, Saucedo, and Pogudin introduce the model and observation used here in equation~(3.3) of version~2 of their preprint~\cite{PSP}. Their Conjecture~1 asks for global structural identifiability for each fixed stage number and known population. Their Definition~1 compares a generic true point with all real competing points of that fixed model. The preceding manuscript~\cite{Fixed} treats that fixed-order question. The present question is whether the observed incidence can also distinguish the stage number, and whether exceptional true points remain within the strictly positive region.

The equal-infectivity corollary below allows every positive true point and every finite positive competing order. Its restriction to positive competitors is essential to the stated scope; it does not replace the all-real-competitor assertion in~\cite{Fixed} or in the definition of~\cite{PSP}. The Darboux approach to epidemiological identifiability already appears in Appendix~B of~\cite{PSP}, for a different model with staged latency. Here the additional issue is to control all positive orbit closures, including those of codimension greater than one.

Fix $N_0>0$. For an integer $n\geq1$, let $J=\sum_{i=1}^n I_i$ and consider
\begin{align}
 \dot S&=-\beta SJ,& \dot I_1&=\beta SJ-\zeta I_1,\notag\\
 \dot I_i&=\zeta(I_{i-1}-I_i)\quad(2\leq i\leq n),&
 \dot R&=\zeta I_n, & y&=\beta SJ.\label{eq:physical}
\end{align}
The population constraint is $S+J+R=N_0$. The unknown rates $\beta,\zeta$ and all initial states $S(0),I_1(0),\ldots,I_n(0),R(0)$ are strictly positive. The observation is the complete analytic germ of $y$ at time zero, in the original physical time. Equality of germs means equality on some neighborhood of zero; the rates, states, and number of stages are not supplied with that germ.

\begin{corollary}[Equal infectivity]\label{cor:equal}
Let a true model~\eqref{eq:physical} have any finite order $n\geq1$ and strictly positive rates and initial states. Let a competing model of any finite order $m\geq1$ have strictly positive rates and initial states and the same known population $N_0$. If their incidence germs in the same physical time agree, then $m=n$, their transmission and progression rates agree, and all corresponding initial states agree.
\end{corollary}

The theorem concerns uniqueness from an exact germ. It supplies no error bound or sampling requirement for finite or noisy observations, and makes no assertion for unknown $N_0$, zero or negative states or rates, or a rescaling of the observed time. The corollary follows from the following weighted result.

\subsection{Unknown stage-dependent infectivity}
Replace $J$ in the incidence and susceptible equations by
\[
 J_w=\sum_{j=1}^n w_jI_j,\qquad w_1=1,\quad w_j>0.
\]
More explicitly, the weighted model is
\begin{align}
 \dot S&=-y,& \dot I_1&=y-\zeta I_1,\notag\\
 \dot I_j&=\zeta(I_{j-1}-I_j)\ (2\leq j\leq n),&
 \dot R&=\zeta I_n,& y&=\beta S J_w.\label{eq:weighted}
\end{align}
Population is still $N_0=S+\sum I_j+R$, using the unweighted sum. The weights $w_2,\ldots,w_n$ are unknown. The normalization $w_1=1$ removes the scale exchange between $\beta$ and all the weights.

\begin{theorem}[Weighted identification]\label{thm:main}
Fix a known $N_0>0$. For any finite $n\geq1$, take strictly positive rates, weights, and all initial states in~\eqref{eq:weighted}, with $w_1=1$ and total population $N_0$. If a competing model of any finite order $m\geq1$, with the same positivity, normalization, and population conditions, has the same incidence germ in the same physical time, then $m=n$, both rates agree, all corresponding weights agree, and every initial state agrees.
\end{theorem}

Linear-chain representations of Erlang and more general dwell-time distributions are established tools; Hurtado and Kirosingh~\cite{HK} develop their construction and generalization. The zero-state input-to-$J_w$ kernel here is
\[
 e^{-\zeta t}\sum_{j=1}^n w_j\frac{(\zeta t)^{j-1}}{(j-1)!}.
\]
Ritschel and Wyller~\cite{RW} estimate mixed-Erlang kernel weights and a common rate through a dynamical maximum-likelihood problem at a prescribed truncation order, with a constant initial history. Our initial infection states are free positive unknowns, so this kernel alone does not describe the full response. Nor is the kernel itself observed: only $y=\beta S J_w$ is observed, with $S(0)$ unknown. Section~\ref{sec:recovery} proves the step from that nonlinear observation to equality of transfer functions before using poles to identify the chain. We make no claim for arbitrary phase-type networks, unequal progression rates, parallel or branching chains, or nonpositive competitors.

\paragraph{Proof and verification scope.}
The theorem is supported by the prose proof below and by exact finite algebra checks described in Appendix~\ref{app:checks}. The accompanying \texttt{VERIFICATION-RECORD.md} describes the scope and dependencies of the AI reviews and finite checks. The article has not received human peer review or formal proof certification; priority has not been established.

\section{Algebraic closures and preliminary exclusions}\label{sec:prelim}
Use the invertible changes
\[
 x=\frac{\beta S}{\zeta},\qquad u_i=\frac{\beta I_i}{\zeta},
 \qquad \tau=\zeta t.
\]
Fix the positive weights and write $D=d/d\tau$ and, in these coordinates, $J=\sum_i w_i u_i$. The polynomial derivation becomes
\begin{equation}\label{eq:D}
 Dx=-xJ,\qquad Du_1=xJ-u_1,\qquad Du_i=u_{i-1}-u_i\quad(2\leq i\leq n).
\end{equation}
All algebraic varieties in the proof are over $\C$. For a positive real initial point $\theta$, let $Z_\theta\subset\A^{n+1}$ be the Zariski closure of its trajectory germ. Evaluation into analytic germs has integral-domain image, so its kernel is a prime ideal. That ideal is $D$-stable, since differentiating a vanishing germ gives another vanishing germ. Thus $Z_\theta$ is irreducible and invariant.

Put
\begin{equation}\label{eq:H}
 A=\sum_{i=1}^n w_i,\qquad \ell_i=\sum_{j=i}^n w_j,\qquad K=\sum_{i=1}^n\ell_i u_i,\qquad z=K+Ax.
\end{equation}
Direct telescoping gives
\begin{equation}\label{eq:basic}
 DK=(Ax-1)J,\qquad Dz=-J=\frac{Dx}{x},\qquad D(z-\log x)=0,
\end{equation}
where the logarithm is used locally away from $x=0$.

\begin{lemma}\label{lem:log}
Let $V$ be an irreducible complex algebraic variety, and let $x,z\in\C(V)$ with $x$ nonconstant. No local branch of $z-\log x$ is constant on a nonempty relative complex analytic open set where it is defined.
\end{lemma}
\begin{proof}
Such an identity gives $dz=dx/x$ as rational differentials: intersect the open set with the smooth locus, and use that a proper algebraic subset cannot contain a relative analytic open set. Extend $x$ to a transcendence basis $x,t_2,\ldots,t_d$ and set $k=\C(t_2,\ldots,t_d)$. The function field $L=\C(V)$ is finite separable over $k(x)$. The unique extension of $\partial_x$ to $L$ satisfies $\partial_xz=1/x$. Its normalized trace
\[
 w=\frac{1}{[L:k(x)]}\Tr_{L/k(x)}z\in k(x)
\]
would have $\partial_xw=1/x$. Trace commutes with the derivation through the embeddings in a normal closure. The derivative of a rational function has zero coefficient of $x^{-1}$ in its Laurent expansion at zero, a contradiction. This reasoning is in the function field and assumes no affine point above $x=0$.
\end{proof}

\begin{proposition}\label{prop:dim}
For every positive $\theta$ and every $n\geq1$, $\dim Z_\theta\geq2$.
\end{proposition}
\begin{proof}
The trajectory is nonconstant, since $Dx=-xJ<0$. If its closure were a curve, its nonzero derivation and~\eqref{eq:basic} would imply $dz=dx/x$ in its function field: the space of rational differentials on a curve has dimension one, and evaluation on a nonzero derivation is injective. The trace argument of Lemma~\ref{lem:log} again gives a contradiction.
\end{proof}

\begin{proposition}[Darboux classification]\label{prop:darboux}
In the equal-infectivity case $w_1=\cdots=w_n=1$, for $n\geq2$, every nonzero polynomial $P$ satisfying $DP=cP$ with polynomial cofactor $c$ has the form $P=A x^k$, where $A\in\C^*$ and $k\geq0$. Its cofactor is $-kJ$.
\end{proposition}
\begin{proof}
The degree of $c$ is at most one. Factor off the largest power of $x$, reducing to $x\nmid P$. Write
\[
 T=-\sum_i u_i\partial_{u_i}+N,\quad
 N=\sum_{i=2}^n u_{i-1}\partial_{u_i},\quad
 D=T-xJ\partial_x+xJ\partial_{u_1},\quad
 P=\sum_{i=0}^s x^iP_i(u).
\]
Here $P_0\ne0$. The equation at $x=0$ is $TP_0=c(0,u)P_0$. Since $T$ preserves degree, $c(0,u)$ is constant. On the homogeneous degree-$r$ space, $T=-r+N$, with $N$ nilpotent. Thus $P_0$ is homogeneous of some degree $r\geq0$, $NP_0=0$, and $c=ax-r$ for a scalar $a$.

The coefficient of $x^i$, for $i\geq1$, gives
\begin{equation}\label{eq:adic}
 (T+r-iJ)P_i=(a-J\partial_{u_1})P_{i-1}.
\end{equation}
For $i=1$ the right side is homogeneous of degree $r$ or is zero. If $P_1\ne0$ has highest degree $d$, the term $-J(P_1)_d$ forces $d\leq r-1$. But its lowest-degree component of degree $e\leq d$ then satisfies $(r-e+N)(P_1)_e=0$, impossible because $r-e>0$ and $N$ is nilpotent. If the right side is zero, the highest-degree term alone gives the contradiction. Hence $P_1=0$ and
\[
 J\partial_{u_1}P_0=aP_0.
\]
Regarding $J=u_1+\cdots+u_n$ as the first coordinate, polynomiality implies $a$ is a nonnegative integer and $P_0=J^a H(u_2,\ldots,u_n)$ for nonzero $H$. If $a>0$, use $NJ=J-u_n$ in $NP_0=0$, divide by $J^{a-1}$, and restrict to $J=0$. This gives $-a u_nH=0$, a contradiction for $n\geq2$. Thus $a=0$. Now $P_0$ is independent of $u_1$; comparing the coefficient of $u_1$ in $NP_0=0$ removes $u_2$, and successively removes every variable. Therefore $P_0$ is constant and $r=0$. Equation~\eqref{eq:adic} and highest-degree comparison give $P_i=0$ for all $i\geq2$. Restoring the factored power of $x$ proves the claim.
\end{proof}

For the equal-infectivity case $n=1$ there are also the factors $u_1^r$; this is a useful control on the restriction $n\geq2$. Proposition~\ref{prop:dim} already gives density for $n=1$. Still in the equal-infectivity setting, for $n=2$, a proper closure would be an invariant hypersurface, whose defining irreducible polynomial is Darboux. Proposition~\ref{prop:darboux} would make that hypersurface $x=0$, which contains no positive point. Higher dimensions require the additional argument below; hypersurface classification alone does not exclude a closure of higher codimension.

\section{The equilibrium line belongs to every positive closure}\label{sec:EL}
Let $L_{\rm eq}=\{(a,0,\ldots,0):a\in\C\}$.
\begin{theorem}[Equilibrium line]\label{thm:EL}
For every $n\geq1$ and every strictly positive real $\theta$, $L_{\rm eq}\subset Z_\theta$.
\end{theorem}
\begin{proof}
The positive orthant is forward invariant. Each infection coordinate solves a linear decay equation with nonnegative input, and $x$ satisfies its exponential equation. Since $D(x+\sum_i u_i)=-u_n\leq0$, the state stays bounded and exists for all forward real time. The conserved quantity $H=z-\log x$ and $K,Ax\geq0$ imply $x\geq\exp(-H(\theta))>0$. Thus $x$ decreases to some $a>0$, and
\[
 \int_0^\infty J(\tau)\,d\tau=\log\frac{x(0)}a<\infty.
\]
The derivative $DJ$ is bounded on the bounded trajectory. Uniform continuity and nonnegativity of $J$ then give $J\to0$: a positive limsup would yield disjoint intervals of fixed width and positive integral. All $u_i$ and $K$ tend to zero because every $w_i$ is strictly positive. If $a\geq1/A$, the identity $DK=(Ax-1)J$ would make $K$ nondecreasing, contrary to $K(0)>0$ and $K\to0$. Hence
\begin{equation}\label{eq:threshold}
 0<a<1/A.
\end{equation}
Analytic continuation of every polynomial vanishing on the trajectory germ shows that the limit $p_a=(a,0,\ldots,0)$ lies in $Z_\theta$.

Write the infection equation as $Du=B_n(x)u$, where $B_n(x)=L-\mathrm{Id}+xe_1w^T$ and $L$ has ones on its subdiagonal and zeros elsewhere. The matrix $B_n(a)$ is Hurwitz. Indeed, if an eigenvalue had $\Re\lambda\geq0$, then $b=\lambda+1$ would have $|b|\geq1$. The eigenvector recurrence gives
\[
 1=a\sum_{j=1}^n\frac{w_j}{b^j},\qquad
 1\leq a\sum_{j=1}^n\frac{w_j}{|b|^j}\leq aA<1,
\]
contrary to~\eqref{eq:threshold}. This includes $n=1$.

Choose a holomorphic logarithm near the positive number $a$. Since $H_x(p_a)=A-1/a\ne0$, the implicit-function theorem gives ambient coordinates $(c,u)$ with $c=H$, $x=\chi(c,u)$, and
\[
 Dc=0,\qquad Du=B_n(\chi(c,u))u.
\]
Let $B_0=B_n(a)$. The positive Hermitian matrix
\[
 M=\int_0^\infty e^{B_0^*\tau}e^{B_0\tau}\,d\tau
\]
exists by the exponential decay of a Hurwitz matrix and satisfies $B_0^*M+MB_0=-\mathrm{Id}$. On a small complex neighborhood of $(H(p_a),0)$, continuity gives
\[
 B_n(\chi(c,u))^*M+M B_n(\chi(c,u))\leq-\tfrac12\mathrm{Id}.
\]
Take a small $c$-disk and an ellipsoid $u^*Mu<\varepsilon^2$ whose closure stays in this chart. Along real time with complex state,
\[
 \frac{d}{d\tau}(u^*Mu)\leq-\tfrac12\|u\|^2
 \leq-\frac{u^*Mu}{2\lambda_{\max}(M)}.
\]
The solution therefore remains in the chart and converges exponentially to
\begin{equation}\label{eq:projection}
 \pi(c,u)=(\chi(c,0),0).
\end{equation}
The map $c\mapsto\chi(c,0)$ is locally injective. This construction takes place in the ambient space and does not require $Z_\theta$ to be smooth.

Suppose $Z_\theta$ did not contain $L_{\rm eq}$. Its intersection with that line would be finite. Shrink the chart so that its limit map meets that intersection only at $p_a$. Every point $q\in Z_\theta$ in the chart remains in $Z_\theta$ under the flow: for finite generators $f_i$ of its $D$-stable ideal, equations $Df_i=\sum_j a_{ij}f_j$ give a homogeneous linear ODE for their values. Algebraic sets are closed in the complex Euclidean topology, so $\pi(q)\in Z_\theta$ and therefore $\pi(q)=p_a$. Local injectivity in~\eqref{eq:projection} makes $H(q)=H(p_a)$ throughout a nonempty relative analytic open subset of $Z_\theta$. But $x$ is nonconstant on $Z_\theta$, since it decreases along the positive trajectory. Lemma~\ref{lem:log} gives a contradiction.
\end{proof}
The real trajectory converges to $p_a$ with $a>0$. The theorem concerns its algebraic closure; it does not assert real convergence to $x=0$.

\section{An algebraic obstruction at the boundary}\label{sec:boundary}
The next theorem needs less than positive weights. In this section alone, fix any complex weights with $w_n\ne0$ and use~\eqref{eq:D} with $J=\sum_iw_i u_i$.
\begin{theorem}[Affine-boundary exclusion]\label{thm:B}
Let $n\geq2$. If $Z\subset\A^{n+1}$ is an irreducible closed invariant variety, $x$ is not identically zero on $Z$, and $2\leq d=\dim Z\leq n$, then $Z\cap\{x=0\}=\varnothing$.
\end{theorem}
The exclusions of dimension one and dimension $n+1$ are necessary: the equilibrium line and the whole ambient space both meet $x=0$.

\subsection{A linear equation over a constant field}
\begin{lemma}\label{lem:obstruction}
Let $F$ be any field of characteristic zero, let $s,t$ be algebraically independent over $F$, and let $\delta=-s\partial_s+\partial_t$ on $F(s,t)$, with $\delta F=0$. For nonzero $P,R\in F[t]$ and nonzero $\lambda\in\Q$, the equations
\begin{align}
 (\delta-\lambda sP)c&=0,\label{eq:hom}\\
 (\delta-sP)c&=R\label{eq:inhom}
\end{align}
have, respectively, no nonzero algebraic solution and no algebraic solution over $F(s,t)$. The one-variable statements on $F(s)$ with $\delta=-s\partial_s$ and nonzero constant $P,R$ hold as well.
\end{lemma}
\begin{proof}
First classify the Darboux factors for $\delta$ on $F[s,t]$. If $\delta f=hf$, total degree makes $h$ constant. Writing $f=\sum_i s^if_i(t)$ gives $f_i'=(h+i)f_i$. A nonzero polynomial coefficient must be constant and have $h=-i$. Thus every nonconstant irreducible Darboux factor is $s$.

If~\eqref{eq:hom} had a nonzero algebraic solution, its norm would give a rational function $v$ with $\delta v/v=e\lambda sP$, where $e$ is the finite extension degree. A numerator or denominator factor other than $s$ gives an uncancelled pole in this logarithmic derivative. Thus $v=a s^b$ for $a\in F^*$ and $b\in\mathbb Z$, and $\delta v/v=-b$, contrary to $e\lambda sP\ne0$.

If~\eqref{eq:inhom} had an algebraic solution, normalized trace would give a rational solution $v$. A pole of order $q$ at an irreducible factor other than $s$ produces a pole of order $q+1$ in $\delta v$, which cannot be cancelled by $-sPv$ or $R$. Hence $v\in F[s,s^{-1},t]$. If its least $s$-exponent $b$ were negative, the coefficient of $s^b$ on the left would be $v_b'-b v_b\ne0$ for nonzero polynomial $v_b(t)$, whereas the right side has no negative $s$-power. Thus $v$ is polynomial. Its greatest $s$-exponent $d\geq0$ then produces the nonzero coefficient $-Pv_d$ at $s^{d+1}$, again impossible. The zero function is not a solution either.

Derivations extend uniquely across separable algebraic extensions and commute with their embeddings, which justifies both norm and trace. The one-variable claim follows by the same argument, or by adjoining an independent $t$ with $\delta t=1$.
\end{proof}

\subsection{The residue field of a boundary component}
Suppose that a variety $Z$ as in Theorem~\ref{thm:B} meets $x=0$, and choose an irreducible component $W$ of that intersection. Since $x$ is a nonzero nonunit on $Z$, the principal-divisor dimension formula gives $\dim W=d-1\geq1$. The local flow preserves the intersection and cannot permute its finitely many irreducible components near the identity, so $W$ is invariant under
\[
 D_0u_1=-u_1,\qquad D_0u_i=u_{i-1}-u_i.
\]
Let $k$ be the first index with $u_k$ not identically zero on $W$. If $k=n$, then $W$ is the last coordinate axis, $d=2$, and its function field is $\C(s)$ with $s=u_n$, $D_0s=-s$, and $J=w_n s$.

For $k<n$, set $q=n-k$ and, on $W$, put
\[
 s=u_k,\qquad t=u_{k+1}/s,\qquad v_i=u_{k+i}/s\quad(0\leq i\leq q).
\]
Then $v_0=1$, $D_0s=-s$, $D_0t=1$, and $D_0v_i=v_{i-1}$ for $i\geq1$. Define
\begin{equation}\label{eq:invariants}
 c_i=\sum_{j=0}^i\frac{(-t)^j}{j!}v_{i-j},\qquad
 v_i=\sum_{a=0}^i\frac{c_a t^{i-a}}{(i-a)!}.
\end{equation}
Here $c_0=1$, $c_1=0$, and $D_0c_i=0$. Set $F=\C(c_2,\ldots,c_q)\subset\C(W)$. These are constants for the boundary derivation, not for the full nonlinear system.

The elements $s,t$ are algebraically independent over $F$. Algebraicity of $t$ over a constant field would force $D_0t=0$. If $s$ were algebraic over $F(t)$, its norm would yield a nonzero rational function $v(t)$ with $v'/v$ equal to minus the extension degree. A rational logarithmic derivative is $O(1/t)$ at infinity, so it cannot be a nonzero constant. Formula~\eqref{eq:invariants} recovers all state coordinates, proving
\begin{equation}\label{eq:residue}
 \C(W)=F(s,t),\qquad D_0=\delta=-s\partial_s+\partial_t,
 \qquad r:=\trdeg_\C F=d-3.
\end{equation}
In particular $\dim W\geq2$. The residue of $J$ has the form $sP(t)$, where
\begin{equation}\label{eq:leadingP}
 P(t)=\sum_{i=0}^q w_{k+i}v_i\in F[t],\qquad
 \deg P=q,\qquad [t^q]P=\frac{w_n}{q!}\ne0.
\end{equation}

\subsection{A valuation chart with moving parameters}\label{sec:valuation}
Choose a prime divisor above $W$ in the normalization of $Z$, and normalize its valuation by $v(x)=1$. Finite normalization and the codimension-one local ring give a discrete valuation; its residue field can be a finite extension of $\C(W)$, which is allowed here. The standard local-ring and valuation facts used at this step are recorded in~\cite{Stacks}. All regular state coordinates have nonnegative valuation, and $s$ is a unit at this generic point.

Choose a transcendence basis $a_1,\ldots,a_r$ of $F/\C$ from among the boundary $c_i$. Lift these elements, $s$, and $t$ to $\C(Z)$ using their actual rational expressions in the state coordinates, and retain the same notation. Their residues are algebraically independent. The elements $x,s,t,a_1,\ldots,a_r$ are therefore algebraically independent as well: after removing the least power of $x$ from a polynomial relation, reduction modulo the valuation would contradict independence of the other residues. There are $r+3=d$ elements, so $\C(Z)$ is finite separable over
\[
 E(x),\qquad E=\C(s,t,a_1,\ldots,a_r).
\]
The valuation is trivial on nonzero elements of $E$ and restricts to the $x$-adic valuation on $E(x)$. Complete at the chosen divisor and enlarge the coefficient field to an algebraic closure $\overline E$. The characteristic-zero Newton--Puiseux theorem embeds the finite extension, compatibly with the chosen valuation, in $\overline E((x^{1/e}))$ for some $e$. We use only this one-variable theorem, with a possibly transcendental coefficient field; see the characteristic-zero statement in~\cite{Poonen}. Both $\C(W)$ and the finite residue extension embed in $\overline E$. Thus the series coefficients are algebraic over $F(s,t)$.

Identify the residue of $E$ with $E$ through the chosen lifts, and embed the residue field over this identification. The embedded boundary field $F$ is algebraic over $\C(a_1,\ldots,a_r)$. Thus the coefficient derivation $\delta=-s\partial_s+\partial_t$ annihilates this copy of $F$, although it need not annihilate general coefficients algebraic over $F(s,t)$. Actual lifts of the invariants need not be $D$-constant.

Each coordinate partial derivation on $E(x)$ extends uniquely over a separable algebraic extension. The termwise coordinate partials on the Puiseux series agree, on the embedded algebraic function field, with these unique extensions. In the series embedding, the actual derivation is
\begin{equation}\label{eq:fullD}
 D=(Dx)\partial_x+(Ds)\partial_s+(Dt)\partial_t
       +\sum_{j=1}^r(Da_j)\partial_{a_j}.
\end{equation}
The exact formulas below will give
\begin{equation}\label{eq:leadingD}
 Dx=-xsP+o(x),\quad Ds=-s+o(1),\quad Dt=1+o(1),
 \quad Da_j=o(1),
\end{equation}
where the notation refers to valuation. The leading coefficient derivation is $\delta=-s\partial_s+\partial_t$. It kills elements algebraic over $\C(a_1,\ldots,a_r)$, in particular the embedded boundary field $F$, but not arbitrary algebraic coefficients over $E$. The actual $Da_j$ need not be zero.

If $h=cx^\nu+$ higher terms is nonzero with $\nu=v(h)>0$, formula~\eqref{eq:fullD} gives
\begin{equation}\label{eq:leadingh}
 [x^\nu]Dh=(\delta-\nu sP)c\ne0.
\end{equation}
Indeed $\partial_x$ lowers exponents by one and the other coefficient partials do not lower them; all parameter-motion terms have strictly higher valuation. The coefficient is nonzero by~\eqref{eq:hom}, since $\nu$ is a positive rational number. Consequently $v(Dh)=v(h)$. This includes ramification, algebraic leading coefficients, and the finite extension of the boundary residue field. For $k=n$, use the basis $x,s$ instead; the same statements hold with $\delta=-s\partial_s$ and $P=w_n$.

\subsection{Components whose first nonzero stage is later than the first}
Let $k\geq2$ and define on $Z$
\[
 s=u_k,\qquad h=u_1/s,\qquad b=u_{k-1}/s,\qquad P^*=J/s.
\]
The residues of $h$ and $b$ are zero, and $P^*$ has the nonzero residue $P$ in~\eqref{eq:leadingP}, or $w_n$ on the last axis. Exact differentiation gives
\begin{equation}\label{eq:later}
 Ds=s(b-1),\qquad Dh=xP^*-hb,\qquad Dx=-xsP^*.
\end{equation}
For $k<n$, one also has $Dt=1-tb$ and $Dv_i=v_{i-1}-bv_i$ for $i\geq1$. The actual lifts of~\eqref{eq:invariants} satisfy
\begin{equation}\label{eq:laterc}
 Dc_i=bG_i,\qquad G_i=t c_{i-1}-c_i+\frac{(-t)^i}{i!}\quad(i\geq2).
\end{equation}
These formulas verify~\eqref{eq:leadingD}, even when $v(b)$ is fractional.

If $h=0$, equation~\eqref{eq:later} would give $xP^*=0$, impossible. Thus $\nu=v(h)>0$. Since $v(hb)>\nu$ (with $v(0)=\infty$), equations~\eqref{eq:leadingh} and~\eqref{eq:later} force $\nu=1$. For $\nu<1$ both right-hand terms have higher valuation, while for $\nu>1$ their sum has valuation one. At valuation one the coefficient equation is
\[
 (\delta-sP)c=P,
\]
contrary to Lemma~\ref{lem:obstruction}. On the terminal axis it is $(\delta-sw_n)c=w_n$, with the same contradiction. This also covers $n=k=2$, where $h=b$.

\subsection{Relations among boundary invariants}
It remains to treat $k=1$. Here $s=u_1$, $t=u_2/u_1$, $P^*=J/s$, and the exact equations are
\[
 Ds=s(xP^*-1),\quad Dt=1-xtP^*,\quad
 Dv_i=v_{i-1}-xP^*v_i\ (i\geq1),\quad Dv_0=0.
\]
Hold a formal variable $T$ constant and work modulo $T^n$. Put
\[
 V(T)=\sum_{i=0}^{n-1}v_iT^i,\qquad
 C(T)=e^{-tT}V(T)=\sum_{i=0}^{n-1}c_iT^i.
\]
Then $DV=TV-xP^*(V-1)$ and
\begin{equation}\label{eq:generating}
 DC(T)=xP^*\bigl[(tT-1)C(T)+e^{-tT}\bigr].
\end{equation}
Coefficient comparison, for arbitrary $n$, yields
\begin{equation}\label{eq:ci}
 Dc_i=xP^*G_i,\qquad
 G_i=t c_{i-1}-c_i+\frac{(-t)^i}{i!},\qquad 2\leq i\leq n-1.
\end{equation}
In particular the chosen parameter lifts satisfy~\eqref{eq:leadingD}.

Let $\mathcal I$ be the kernel of the evaluation map
$\C[C_2,\ldots,C_{n-1}]\longrightarrow F$ that sends $C_i$ to the boundary residue of $c_i$. It is a proper prime ideal. If $n=2$,~\eqref{eq:residue} already gives $\dim W\geq2$, incompatible with $d\leq2$. If $n\geq3$, then
\[
 \trdeg_\C F=d-3\leq n-3<n-2,
\]
so $\mathcal I$ must be nonzero.

Take $f\in\mathcal I$ and its actual lift $h=f(c_2,\ldots,c_{n-1})$. Its residue is zero, and~\eqref{eq:ci} gives
\begin{equation}\label{eq:relation}
 Dh=xP^*R(t,c),\qquad R=\sum_{i=2}^{n-1}G_i\partial_i f.
\end{equation}
Suppose the residue $R_0\in F[t]$ were nonzero. If $h=0$, equation~\eqref{eq:relation} is already a contradiction. Otherwise~\eqref{eq:leadingh} forces $v(h)=1$ and gives
\[
 (\delta-sP)c=PR_0,
\]
contrary to~\eqref{eq:inhom}. Thus $R_0=0$ for every $f\in\mathcal I$.

The variable $t$ is transcendental over $F$. For each $i\geq2$, the coefficient of $t^i$ in $R_0$ is exactly $(-1)^i(\partial_i f)(c)/i!$: all terms $t c_{j-1}-c_j$ have degree at most one, and the remaining term of $G_j$ has degree exactly $j$. Hence $\partial_i f\in\mathcal I$ for every $i$. In characteristic zero a proper polynomial ideal stable under all coordinate partial derivatives must be zero. To see this, choose a nonzero member of least total degree; if nonconstant, it has a nonzero partial derivative of lower degree, and if constant, it cannot belong to a proper ideal. This contradicts $\mathcal I\ne0$.

All possible $W$ have been excluded, proving Theorem~\ref{thm:B}. The case of no relations would give $\dim W=2+(n-2)=n$ and $\dim Z=n+1$, exactly the ambient-space case deliberately excluded from the theorem.

\section{Density of every positive trajectory}
\begin{theorem}\label{thm:density}
For every finite $n\geq1$, every fixed positive weight vector with $w_1=1$, and every positive real initial point of~\eqref{eq:D}, the trajectory is Zariski dense in $\A^{n+1}$.
\end{theorem}
\begin{proof}
Theorem~\ref{thm:EL} puts the equilibrium line and its origin in $Z_\theta$. Since the closure is irreducible and also contains a positive point outside that line, its dimension is at least two, consistently with Proposition~\ref{prop:dim}. For $n=1$ this is already full dimension. For $n\geq2$, a proper closure would have dimension between two and $n$, and $x$ would not vanish identically on it. Theorem~\ref{thm:B} would make its intersection with $x=0$ empty, contradicting the origin supplied by Theorem~\ref{thm:EL}.
\end{proof}
The density assertion concerns the reduced state space after eliminating $R$. With $R$ included, the conserved population imposes its affine constraint. No density in an unconstrained space including $R$ is asserted.

\section{Recovering the hidden chain in physical time}\label{sec:recovery}
Return to~\eqref{eq:weighted}, and write $D=d/dt$ throughout this section. The invertible state scaling and the nonzero time scaling used earlier transfer density to the true trajectory in $\C[S,I_1,\ldots,I_n]$ at its fixed rates and weights. Both observations in the argument that follows use the same original physical time.

For the linear infection chain define
\begin{equation}\label{eq:Bn}
 B_{n,w}(X,\zeta)=\sum_{j=1}^n w_j\zeta^{j-1}(X+\zeta)^{n-j}.
\end{equation}
The identities $(D+\zeta)I_1=y$ and $(D+\zeta)I_j=\zeta I_{j-1}$ imply
\begin{equation}\label{eq:chain}
 (D+\zeta)^nJ_w=B_{n,w}(D,\zeta)y.
\end{equation}
This holds with arbitrary initial infection states; no zero-state assumption is used.

\subsection{The susceptible offset vanishes}
Let a positive competing model have order $m$, rates $b,\eta$, and weights $v_1=1,v_2,\ldots,v_m>0$. Equality of incidence germs gives $S^*=S+\Delta$ for a constant $\Delta$. Set $q=S+\Delta$, so $Dq=-y$ and the competing pressure is $J^*=y/(bq)$. Division is valid near the initial time because both susceptible states are positive.

Define polynomials, using differentiation along the true model, by
\begin{equation}\label{eq:Pj}
 P_0=y,\qquad P_{j+1}=qDP_j+(j+1)yP_j.
\end{equation}
The quotient rule gives $D^j(y/q)=P_j/q^{j+1}$. Substitute into the competing chain identity and clear denominators to obtain the state polynomial
\begin{equation}\label{eq:Fm}
 F_m=\sum_{j=0}^m\binom mj\eta^{m-j}q^{m-j}P_j
       -bq^{m+1}B_{m,v}(D,\eta)y.
\end{equation}
It vanishes on the true trajectory. Theorem~\ref{thm:density} therefore makes it the zero polynomial. Reducing~\eqref{eq:Pj} modulo $q$ gives
\begin{equation}\label{eq:remainder}
 F_m\bmod q=m!y^{m+1}.
\end{equation}
If $\Delta\ne0$, substitution of $S=-\Delta$ gives $y=-\beta\Delta\sum w_j I_j$, a nonzero state polynomial. Equation~\eqref{eq:remainder} is then impossible. Hence $\Delta=0$ and $S^*=S$. The negative substitution is a test of a polynomial identity, not the introduction of negative competing physical states.

\subsection{Rates, order, and weights}
Now $J^*=\beta J_w/b$. Eliminating $J_w$ between the true and competing versions of~\eqref{eq:chain} gives $T(D)y=0$, where
\begin{equation}\label{eq:T}
 T(X)=\beta(X+\eta)^m B_{n,w}(X,\zeta)
       -b(X+\zeta)^n B_{m,v}(X,\eta).
\end{equation}
We show that $T$ is the zero polynomial in $X$.

The highest homogeneous part of the physical vector field, acting on functions of $S,J_w$, is
\[
 D_2=\beta S J_w(-\partial_S+\partial_{J_w}),
\]
because $w_1=1$. The highest degree of $D^k y$ is $k+2$, with nonzero coefficient $(-1)^k\beta^{k+1}$ on $S J_w^{k+1}$. Indeed, to obtain that monomial from $D_2^k(\beta S J_w)$ one must choose the $-\partial_S$ term at every step; choosing the other term increases the exponent of $S$, which subsequent steps cannot decrease. The linear part of the vector field cannot contribute to this highest degree. Since $S,J_w$ are algebraically independent, these leading polynomials are nonzero. Distinct $k$ give distinct degrees, so $y,Dy,\ldots$ are linearly independent over $\C$ as state polynomials. Density applied to $T(D)y=0$ proves $T=0$.

The leading coefficient of $B_{n,w}$ is $w_1=1$ and its degree is $n-1$. The highest coefficient of~\eqref{eq:T} is therefore $\beta-b$, giving $b=\beta$. It follows that
\begin{equation}\label{eq:transfer}
 \frac{B_{n,w}(X,\zeta)}{(X+\zeta)^n}
   =\frac{B_{m,v}(X,\eta)}{(X+\eta)^m}.
\end{equation}
There is no pole cancellation, since
\[
 B_{n,w}(-\zeta,\zeta)=w_n\zeta^{n-1}\ne0
\]
and the same holds for the competitor. The unique pole and its order give $\eta=\zeta$ and $m=n$. Equality of numerators in the basis $(X+\zeta)^{n-j}$, together with $\zeta\ne0$, gives $v_j=w_j$ for every $j$.

Thus the transfer identity is derived from the nonlinear observation, through density and the independence of the derivatives of $y$. It is not assumed to be directly measured. Since the argument applies to each arbitrary finite $m$, it needs no common finite bound on competing orders and no additional generic restriction.

\subsection{Every initial state}
The common values of $S$ and $\beta$ give the common germ $J_w=y/(\beta S)$. Let $I$ be the infection column vector, $e=e_1$, $L$ the subdiagonal shift, and $A_\zeta=\zeta(L-\mathrm{Id})$. Then $DI=A_\zeta I+ey$. Repeated differentiation gives
\begin{equation}\label{eq:observation}
 D^rJ_w=w^TA_\zeta^r I+
 \sum_{h=0}^{r-1}(w^TA_\zeta^{r-1-h}e)D^h y,
 \qquad r\geq0,
\end{equation}
with an empty sum when $r=0$. For $r=0,\ldots,n-1$, subtract the known input terms. The coefficient matrix $\mathcal O$ has rows $w^TA_\zeta^r$. Unit lower-triangular binomial row operations change these rows to $w^T(A_\zeta+\zeta\mathrm{Id})^r=\zeta^r w^TL^r$. This matrix is triangular about its opposite diagonal, with opposite-diagonal entries $\zeta^r w_n$, so
\begin{equation}\label{eq:obsdet}
 \det\mathcal O=(-1)^{n(n-1)/2}\zeta^{n(n-1)/2}w_n^n\ne0.
\end{equation}
Equation~\eqref{eq:observation} uniquely recovers the full infection vector, in particular its initial value. Finally
\[
 R(0)=N_0-S(0)-\sum_{j=1}^n I_j(0).
\]
The population uses the unweighted infection sum, not $J_w$. This proves Theorem~\ref{thm:main}.

For $w_j=1$, all the weighted formulas specialize to the equal-infectivity model: $A=n$, $\ell_j=n-j+1$, and
\begin{equation}\label{eq:equalB}
 B_{n,1}(X,\zeta)
 =\sum_{a=0}^{n-1}\binom n{a+1}\zeta^{n-1-a}X^a
 =\frac{(X+\zeta)^n-\zeta^n}{X}.
\end{equation}
The last expression denotes a polynomial identity, including at $X=0$. Restricting both models to these weights proves Corollary~\ref{cor:equal}. The weighted theorem also distinguishes an equal-infectivity true model from positive normalized weighted competitors.

\section{Sharp finite jets against one infectious compartment}\label{sec:finitejet}
The identification theorem uses a complete analytic germ and allows every
finite competing order. For a fixed pair of small orders, a finite list of
exact derivatives already determines the output germ. We give sharp bounds
in a larger compartment class, and then restrict the upper bounds to the
weighted chains of Theorem~\ref{thm:main}.

\subsection{Model class and quantifiers}
Fix a real population $N>0$. Let $\mathcal M_d(N)$ consist of systems
\begin{equation}\label{eq:fjmodel}
 \dot S=-y,\qquad \dot X=KX+e_1y,\qquad
 \dot R=-\mathbf1^TKX,\qquad y=\beta SCX,
\end{equation}
where $X\in\R^d$, $K$ is a real Hurwitz Metzler matrix with
$\mathbf1^TK\leq0$ entrywise, $\beta>0$, $C$ is a strictly positive row,
and $Ce_1=1$. All initial components, including $S(0)$ and $R(0)$,
are strictly positive and $S(0)+\mathbf1^TX(0)+R(0)=N$.
Every parameter and initial state may differ between the two models.
The index $d$ counts infectious states; $S$ and $R$ are additional states.
The polynomial equations have analytic local solutions. Outputs are
compared at zero in the same physical time.

Write $j^b y=(y(0),y'(0),\ldots,y^{(b)}(0))$. Define $B_N(1,m)$ as the
least nonnegative integer $b$ such that
\begin{equation}\label{eq:fjquantifier}
 \text{for every }M_L\in\mathcal M_1(N),\ M_R\in\mathcal M_m(N),\quad
 j^by_L=j^by_R\ \Longrightarrow\ [y_L]_0=[y_R]_0.
\end{equation}
Here $[y]_0$ denotes an analytic germ, and $b$ is the highest derivative
order, so the jet has $b+1$ entries. No controllability or observability
assumption enters this definition. A realization is called minimal when
$(K,e_1)$ is controllable and $(C,K)$ is observable. This separate use of
``minimal'' concerns linear realization dimension, not the least $b$.

\begin{theorem}[Sharp output-germ bounds]\label{thm:fjsharp}
For every real $N>0$,
\[
 B_N(1,1)=4,\qquad B_N(1,2)=8,\qquad
 B_N(1,3)=11,\qquad B_N(1,4)=15.
\]
The upper bounds hold for all pairs in the indicated classes, including
nonminimal realizations. For each entry $b$, there is a pair with both
realizations minimal for which $j^{b-1}y_L=j^{b-1}y_R$ and
$y_L^{(b)}(0)\ne y_R^{(b)}(0)$.
\end{theorem}

The proof below separates analytic reductions from finite algebraic
certificates. The certificate dependencies and replay entry points are
listed in \texttt{FINITE-JET-REPRODUCE.md}; the numerical values alone are
not a proof of the universal upper bounds.

\subsection{Two reductions}
\begin{lemma}[Changing population without changing time]\label{lem:fjscale}
If $a>0$, the map
\[
 (S,X,R,\beta,K,C,y)\longmapsto
 (aS,aX,aR,\beta/a,K,C,ay)
\]
is a bijection from $\mathcal M_d(N)$ to $\mathcal M_d(aN)$.
It preserves minimality and the first derivative order at which two
outputs differ, when applied to both models with the same $a$.
\end{lemma}
\begin{proof}
The observation transforms as $(\beta/a)(aS)C(aX)=ay$.
Each differential equation and the population sum is multiplied by $a$.
Positivity is preserved. The matrices $K,C$, and $e_1$ are unchanged,
so their sign, stability, normalization, and rank conditions are preserved.
For every $j\geq0$, $\widehat y^{(j)}=a y^{(j)}$ in the original time.
The inverse uses $1/a$.
\end{proof}

For the next reduction let $\partial=d/dt$, let $T$ denote the competing
susceptible state, and put $h=T(0)-S(0)$, $q=S+h$, $z=y/q$ on the left
trajectory. Write
\[
 D(s)=\det(sI-K_R),\qquad P(s)=C_R\operatorname{adj}(sI-K_R)e_1.
\]
Then $D$ is monic of degree $m$ and $P$ is monic of degree $m-1$.
Cayley--Hamilton gives
\[
 D(\partial)z_R=\beta_R P(\partial)y_R,
 \qquad z_R=\beta_R C_RX_R=y_R/T,
\]
with arbitrary initial infection states; no cancellation of $D$ is made.

\begin{lemma}[Contact and feedback]\label{lem:fjcontact}
If $b\geq m-1$ and $j^by=j^by_R$, the residual
\[
 \mathcal R=D(\partial)z-\beta_R P(\partial)y
\]
has order at least $b+1-m$. Conversely, if $q(0)=T(0)$ and
$z^{(j)}(0)=z_R^{(j)}(0)$ for $0\leq j<m$, and, for a nonnegative integer $\ell$,
$\mathcal R^{(j)}(0)=0$ for $0\leq j\leq\ell$, then the outputs agree
through order $m+\ell$. Their next difference is
\begin{equation}\label{eq:fjnext}
 y^{(m+\ell+1)}(0)-y_R^{(m+\ell+1)}(0)
   =q(0)\mathcal R^{(\ell+1)}(0).
\end{equation}
If $\mathcal R$ vanishes identically and these initial values agree, the
output germs agree.
\end{lemma}
\begin{proof}
The forward implication follows from $y-y_R=O(t^{b+1})$,
$q-T=O(t^{b+2})$, and $z-z_R=O(t^{b+1})$, since $q(0)=T(0)>0$.
Differentiating at most $m$ times gives the asserted residual order.
For the converse, use $q'=-qz$ and the monic equation for $z^{(m)}$.
Its right side uses only $z$ derivatives through $m-1$ and $y=qz$
derivatives through $m-1$. Induction on the residual derivative index
therefore gives equality of $z$ and $y$ through $m+\ell$, and of $q$
through $m+\ell+1$. Subtract the next equations to obtain
\eqref{eq:fjnext}. With identically zero residual the variables
$(q,z,z',\ldots,z^{(m-1)})$ solve the same polynomial differential
system with the same initial value; local uniqueness proves the last claim.
\end{proof}

\subsection{Upper bounds and their algebraic certificates}
For $m=1$, let $p_j=\partial^j(y_L-y_R)$ in the polynomial ring in both
sets of rates and states, with differentiation along the product system.
An exact polynomial certificate gives
\begin{equation}\label{eq:fjideal}
 p_5=\sum_{j=0}^4 a_jp_j.
\end{equation}
Its five coefficient polynomials have respectively $593,311,127,29,8$
terms. Checking every coefficient establishes an identity for all parameter
values. Along any solution, $(p_0,\ldots,p_4)$ satisfies a homogeneous
linear system with continuous coefficients. Its zero initial value forces
it to vanish, proving the upper bound four.

For $m=2,3,4$, a common positive change of state and time units makes
$\beta_L=S_L(0)=1$, $X_L(0)=r>0$, and $K_L=k<0$.
Explicitly, with $s_0=S_L(0)$ and $c=\beta_Ls_0$, divide all states
by $s_0$ and use $\tau=ct$ on both sides. Then
$\widetilde K=K/c$, $\widetilde\beta=\beta/\beta_L$, and
$\widetilde y=y/(s_0c)$. A $j$th derivative difference is multiplied by
$1/(s_0c^{j+1})$, a nonzero number. This common normalization preserves
contact order; it does not allow a separate time change for each model.
The left first integral is
\[
 X=r+1-S-k\log S.
\]
In these normalized coordinates $h=T(0)-1>-1$ and $q=S+h$.
The polynomial residual $F=q^{m+1}\mathcal R/y$ is an exponential
polynomial in $x=\log S$. Its order is the same as that of
$\mathcal R$, since $q(0),y(0)>0$ and $x'(0)=-r\ne0$.

\begin{lemma}[Certified small-order elimination]\label{lem:fjelim}
Under the preceding positivity conditions, the residuals for
$m=2,3,4$ cannot have, respectively, orders at least $7,9,12$ when
$h\ne0$. When $h=0$, those orders force $F$ to vanish identically.
No minimality assumption on the right realization is needed.
\end{lemma}
\begin{proof}[Proof by algebraic certificates]
For $h=0$ the residual spaces have dimensions $4,7,11$.
They are spanned by exponential monomials with nonzero confluent
Vandermonde determinant at the initial point, and are invariant under
$d/dx$. Vanishing to at least the stated orders therefore gives $F=0$.
For $m=3,4$ the computed initial jet determinants in the chosen bases are
$192$ and $1146617856$. Lemma~\ref{lem:fjcontact} returns to the germ.

For $m=2$, $F$ lies in the eleven-dimensional space with basis
\[
 1,x,x^2,e^x,xe^x,x^2e^x,e^{2x},xe^{2x},e^{3x},xe^{3x},e^{4x}.
\]
The order-seven subspace has dimension four. Write
$F=\lambda p(d/dx)W_0$, where $W_0$ is the solution of order ten,
$p(0)=1$, and $\deg p\leq3$, and put $a=p(1)$ and $b=p(2)$.
Coefficient comparison
reduces the general branch to two polynomial equations in two variables.
Their resultant has the factors $(12a+1)^7$, a cubic, and a degree-twelve
polynomial. The first factor gives $h=-1$; on the cubic the original
polynomials drop degree and a B\'ezout identity excludes a finite solution.
The degree-twelve factor has exactly four real roots. Their unique lifts
violate, in order, positivity of the infection-rate ratio, positivity of
the coefficient $d_1$ of $D$, positivity of $d_1$, and $r>0$.
The ratio-one branch must be checked in all coefficients: its remaining
residual is $3791k^2S/972-10013k^2/8748$, which is nonzero.
The zero denominators are thus handled before any division is used.

For $m=3$, the nine residual equations have a common kernel with last
coordinate one. A zero leading six-row pivot forces four determinant
conditions, already excluded by an all-rank certificate. On the remaining
branch, exact substitution and projection certificates give a squarefree
degree-289 polynomial in $h$ with unique finite lifts in the other two
coordinates. They cover all candidates, including the component
intersections. There are exactly 28 roots with $h>-1$. Thirteen violate
$r>0$, five violate $k<0$, two violate the positive rate ratio, one violates
$d_2>0$, six violate $d_2^2-3d_1\geq0$, and the last violates positivity
of the constant coefficient of the free-response Laplace numerator.
These are necessary conditions for every positive Hurwitz Metzler
realization: in particular that last coefficient equals
$d_0C(-K)^{-1}X(0)>0$. The trace inequality follows by writing
$d_2=-\operatorname{tr}K$ and
$d_1=((\operatorname{tr}K)^2-\operatorname{tr}(K^2))/2$ and using
$K_{ij}K_{ji}\geq0$. Failure to find a particular realization would
not suffice for this exclusion.

For $m=4$, twelve contact rows have a common nonzero kernel.
The determinant of the first nine rows and its first three derivatives
vanish at all ranks. After removing units and putting $u=r/(1+h)$,
this gives $F_4=WF_4=W^2F_4=W^3F_4=0$, where
\[
 W=-h(1+h)\partial_h-k(1+h)\partial_k-(u+1+k)\partial_u.
\]
Here $u$ is an auxiliary coordinate, not physical time.
The normalized $F_4$ has 2754 terms and multidegree $(25,16,12)$.
Integer resultant identities split the projection into factors of
$k$-degrees 82 and 99. Exact divisibility, degree preservation modulo a
good prime, and coprimality of the residual projections cover all possible
$h$ by $h=0,-1$ and seven factors of total squarefree degree 30.
The corresponding number-field fibers have a unique $k$ and polynomial
gcd $u$ in the last coordinate. Thus $u=r=0$ on every remaining fiber,
contrary to positivity. The 16 real roots of the seven factors include
10 with $h>-1$; all are excluded by the same fiber result.

The integer resultant identities are certified by coefficient bounds
$B,C$ for the raw resultant and the candidate after multiplication back,
a prime product $M>B+C$, and congruence of all raw coefficients.
The divisibility factors are proved over the rationals, including the
required $H_{50}\operatorname{rad}(H_{50})$ factor, where $H_{50}$
is the degree-fifty boundary polynomial with squarefree part of degree
thirty, and the local
$(h+1)$ valuations 479 and 500. Coprimality over a finite field is used
only after these divisibilities and preservation of the degree of one
projection have been checked. Neither a nonzero leading contact pivot
nor a nonzero leading coefficient in $u$ is assumed. These certificate
obligations, not modular root counts alone, justify the global exclusion.
\end{proof}

Lemma~\ref{lem:fjcontact} converts output orders $8,11,15$ to residual
orders $7,9,12$. Lemma~\ref{lem:fjelim} and the polynomial identity
\eqref{eq:fjideal} prove the four upper bounds.

\subsection{Witnesses and sharpness}
For $m=1$, the two tuples $(\beta,k,S(0),X(0),R(0))$ are
\[
 (1,-1,1,2,1),\qquad (3/2,-4,2,2/3,4/3).
\]
Both have population four. Their derivatives through order four are
$(2,-4,4,24,-192)$ and $(2,-4,4,24,-208)$.
For the other cases the witnesses are specified by exact algebraic
numbers or unique roots in rational boxes, not rounded displayed values.
The certificates verify positivity, population, and controllability and
observability in addition to contact.

\begin{center}
\begin{tabular}{@{}rrrrl@{}}\toprule
$m$&base $N$&equal through&first difference&certificate type\\\midrule
1&4&3&4&rational pair and polynomial identity\\
2&12&7&8&cubic number field\\
3&100&10&11&isolated root of eight equations\\
4&100&14&15&isolated root of three equations\\\bottomrule
\end{tabular}
\end{center}
For $m=2$ the real embedding is the unique real root of
\[
 110119472\kappa^3+2679399000\kappa^2+20959796655\kappa
   +51676380000=0.
\]
Exact quotient-field recursion gives equality through order seven and
$55078505<y_R^{(8)}(0)-y_L^{(8)}(0)<55078506$.
For $m=3$ an interval contraction proves the existence of the common
root of eight residual equations. Positive realization and initial
moment matching, followed by Lemma~\ref{lem:fjcontact}, give equality
through order ten and a strictly positive eleventh difference.

For $m=4$, solve the first eight residual equations for
$(d_3,d_2,d_1,d_0,g,b_2,b_1,b_0)$ as functions of $(h,k,r)$.
All elimination pivots exclude zero on the rational box.
The three remaining equations $\mathcal R_8=\mathcal R_9=
\mathcal R_{10}=0$ have a unique root there: an outward-rounded
integer interval calculation maps the box strictly into itself with
contraction bound less than $5.2\cdot10^{-25}$.
At this same root, $D$ has four distinct negative real roots. They define
a lower bidiagonal $K$ with subdiagonal entries one. The first four
Markov coefficients define $C$, and the first four free-response moments
define $X(0)$. The certificates give $C>0$, $Ce_1=1$, $X(0)>0$,
negative column sums, and nonzero controllability and observability
determinants. Both removed initial states are positive at population 100.
Lemma~\ref{lem:fjcontact} proves exact equality through order fourteen;
the certified next difference is
\[
 128995100000000000<y_L^{(15)}(0)-y_R^{(15)}(0)
       <128995200000000000.
\]
Zero containment of interval differences is only a consistency check.
Exact equality follows from the root equations and the feedback induction.
A single such root suffices; a classification of every solution of the
eleven-contact equations is unnecessary for the lower bound.

Finally Lemma~\ref{lem:fjscale}, with $a=N/N_{\rm base}$, transports
each witness to every real $N>0$. Together with the upper bounds this
proves Theorem~\ref{thm:fjsharp}.

\subsection{What the bounds imply for weighted chains}
\begin{corollary}[A bounded-order comparison within the weighted class]\label{cor:fjW}
Let the true model in Theorem~\ref{thm:main} have one infectious stage,
and restrict the competing stage number to $m\in\{1,2,3,4\}$.
For a fixed $m$, equality of the output derivatives through
$b_m=4,8,11,15$, respectively, implies $m=1$ and equality of both rates,
all weights, and every initial state. In particular, derivatives through
order fifteen suffice uniformly over this restricted set of competitors.
\end{corollary}
\begin{proof}
For a weighted chain, $K=\zeta(L-I)$ and $C=(w_1,\ldots,w_m)$.
It is Hurwitz Metzler, its column sums are zero except for the last
sum $-\zeta$, and $Ce_1=w_1=1$. Thus both models belong to the classes
of Theorem~\ref{thm:fjsharp}. The appropriate jet equality first gives
germ equality. Theorem~\ref{thm:main} then gives the conclusion, with
$R(0)$ recovered from the common population and the unweighted state sum.
\end{proof}

For the larger class~\eqref{eq:fjmodel}, equality of germs alone does not
assert recovery of every entry of $K,C$ or every physical state. Conversely,
the lower-bound witnesses above need not be common-rate weighted chains.
Their sharpness therefore does not establish sharpness within that smaller
class for $m=2,3,4$. The rational $m=1$ witness is already in the weighted
class and gives sharpness four for that comparison. The table also supplies
no bound for all pairs whose two orders are at most four, nor a uniform
finite bound over unrestricted competing orders. These are statements
about exact derivatives, not a sampling prescription or a noise estimate.

\subsection{Relation to established methods}
Finite Taylor reduction for generic structural identifiability is
established in Hong, Ovchinnikov, Pogudin, and Yap~\cite{FJHOPY}.
General multiplicity estimates along polynomial vector fields, such as
those of Gabrielov~\cite{FJGAB} and Binyamini~\cite{FJBIN}, give uniform
finite contact bounds after passing to the product of two fixed-order
systems. Positivity makes that vector field nonsingular at the initial
point because $S'=-y<0$. Thus existence of a finite bound and the use
of polynomial elimination are not new general methods here.
Ovchinnikov, Pogudin, and Thompson~\cite{FJOPT} also distinguish
input-output identifiability from single-experiment identifiability;
an input-output equation by itself is not enough to identify its
coefficients from one trajectory.
The contribution asserted in this section is the sharp list for the
specified feedback class and its restricted implication for weighted
chains. The existing literature comparisons do not establish priority
for the new sharp list; model-specific improvements and a new general
method are separate questions.

\paragraph{Use of AI.}
Codex also assisted with the finite-jet synthesis and certificate replay. AI assistance was used in developing the arguments, writing and running exact algebra checks, staged proof review, and preparation of the manuscript. The author retains responsibility for the mathematical claims, attribution, and any subsequent dissemination.

\appendix
\section{Exact finite checks}\label{app:checks}
The five original checks described in this appendix use integer and rational arithmetic from the standard library. The separate finite-jet certificate replays and the additional global-elimination dependencies are described in \texttt{FINITE-JET-REPRODUCE.md}. Its finite checks cover the displayed identities and representative specializations; the all-order results rest on the proofs above. The runner writes to a selected output directory and does not need the source repository history, network access, or earlier work directories.

For equal infectivity, it checks $253$ small Darboux coefficient matrices: $125$ for $n=2$ at total degree at most four, $64$ for $n=3$ at degree at most three, and $64$ at $n=1$ as a control with the additional factors $u_1^r$. It also checks the contracting equilibrium calculation at $n=4$, including a threshold control, and $41$ groups of free-parameter Laurent polynomial identities for the $n=5$ boundary argument. In the latter, the $t^2,t^3,t^4$ coefficient matrix separating the invariant partial derivatives has determinant $-1/288$; the coordinate Jacobian is $1/u_1^4$. Free invariant parameters describe the whole boundary space, a necessary control on the dimension restriction.

For the weighted model, separate finite checks include $25$ reconstruction identities for $n=2,\ldots,6$, $27$ characteristic-polynomial evaluations for $n=1,\ldots,6$, six checks each of~\eqref{eq:obsdet}, weight recovery, and the weighted conservation identity, and seven denominator-remainder checks through order six. Twelve tests recover arbitrary rational infection states after subtracting the input terms in~\eqref{eq:observation}, using both nonuniform and equal weights for $n=1,\ldots,6$. Another $21$ coefficient comparisons verify~\eqref{eq:equalB}. The input derivatives in the state-recovery tests are synthetic values for an exact linear identity, not simulated epidemic observations.

The reproduction command is
\begin{center}\texttt{python3 -B reproduce.py -{}-output build/checks}\end{center}
from the accompanying package directory. It checks file integrity, runs the included bounded checks, and compares their scientific output with frozen baselines while recording runtime and peak memory separately. The package also includes negative controls for a false moving-invariant identity and for tampering with a checked input. A passing finite run checks the listed algebra and packaging; it does not machine-certify complex convergence, normalization, or the valuation argument.

\section{Further fixed-order bounds}\label{sec:fjextension}
For the same positive compartment class and the same physical time used
in Section~7, the following bounds hold for every common real population
$N>0$:
\[
18\le B_{\min}(1,5)\le42,\qquad
10\le B_{\min}(2,2)\le75,\qquad B_{\min}(2,3)\le278.
\]
On the $(2,3)$ branch where the two initial susceptible states coincide,
order $30$ suffices. These statements are bounds, not optimal values.
The previously proposed lower bound $10$ for $(2,3)$ has not been
rechecked in this extension and is not adopted here as a verified bound.
The upper bounds cover nonminimal realizations; the two adopted lower
witnesses are minimal on both sides. The object recovered is the output
germ, not an arbitrary physical realization in the larger class.

The lower bound $18$ uses a certified root of three contact equations,
a positive five-state realization with a cyclic three-state block,
exact agreement through derivative $17$, and a nonzero derivative-$18$
difference. The upper bound $42$ returns to the original five-order
Cayley--Hamilton residual: contact gives order at least $38$ in a
$50$-dimensional exponential-polynomial space. Degree-$11$
interpolation leaves two free coefficients. The next two logarithmic
layers give necessary equations without discarding degenerate pivots;
three certified integer resultants have common gcd $1$ over $\mathbb Q$.
The remaining equal-susceptibility branch yields a closed analytic ODE.
The complete argument, three certificates and both exact checkers are
bundled under \texttt{extensions/u42/}.

For $(2,2)$, the lower witness agrees through order $9$ and differs at
order $10$. The upper bound uses a universal residual resultant in an
$82$-dimensional space, degree-$8$ interpolation, and exact elimination
with real-root exclusion. For $(2,3)$, the nonzero residual resultant
has a $275$-dimensional support envelope after universal coefficient
cancellations. The three-order residual transfers this bound to order
$278$; the equal-susceptibility branch is treated separately at order
$30$. Full proof excerpts, the universal-identity calculations and
bounded reproduction commands are in \texttt{extensions/README.md}.
These are adopted results with their supporting proofs, rather than new
claims of optimality or a new general method.

\paragraph{Verification scope.}
The bounds combine prose analytic arguments, exact algebraic
certificates and staged, nonblind checks. Finite replay certifies the
stated algebraic and witness obligations; the analytic implications
remain prose proofs. The per-claim verification scope, prior exposure,
and procedural qualifications are recorded in
\texttt{extensions/README.md}.


\begin{thebibliography}{9}\small\raggedright
\bibitem{PSP} B. Pant, O. Saucedo, and G. Pogudin.
\emph{Uncovering identifiability of epidemiological models: basic reproduction number and complementary data streams}.
medRxiv, version 2, 25 January 2026.
\href{https://doi.org/10.64898/2026.01.16.26344284}{doi:10.64898/2026.01.16.26344284}.
Definition 1, equation (3.3), Conjecture 1, and Appendix B.
\bibitem{Fixed} T. Takimoto.
\emph{Global identifiability for the staged-infection SIR model}.
Version 1.0.0, 3 November 2026.
\href{https://doi.org/10.5281/zenodo.23900995}{doi:10.5281/zenodo.23900995}.
Same-author preceding work on each fixed known order.
\bibitem{HK} P. J. Hurtado and A. S. Kirosingh.
\emph{Generalizations of the `Linear Chain Trick': Incorporating more flexible dwell time distributions into mean field ODE models}.
\href{https://arxiv.org/abs/1808.07571v1}{arXiv:1808.07571v1} (2018).
\bibitem{RW} T. K. S. Ritschel and J. Wyller.
\emph{An algorithm for distributed time delay identification without a priori knowledge of the kernel}.
Automatica 178 (2025), 112382.
\href{https://doi.org/10.1016/j.automatica.2025.112382}{doi:10.1016/j.automatica.2025.112382}.
\bibitem{Poonen} B. Poonen.
\emph{Maximally complete fields}.
L'Enseignement Math\'ematique 39 (1993), 87--106.
\href{https://math.mit.edu/~poonen/papers/amsval.pdf}{Author manuscript}, Introduction (characteristic-zero Puiseux statement).
\bibitem{Stacks} The Stacks Project Authors.
\emph{The Stacks Project}.
\href{https://stacks.math.columbia.edu/tag/00PD}{Tag 00PD} (normal local rings of dimension one) and
\href{https://stacks.math.columbia.edu/tag/0ASF}{Tag 0ASF} (finite extensions of discrete valuation fields).
\bibitem{FJHOPY} H. Hong, A. Ovchinnikov, G. Pogudin, and C. Yap.
\emph{Global Identifiability of Differential Models}.
Communications on Pure and Applied Mathematics (2020).
\href{https://arxiv.org/abs/1801.08112v5}{arXiv:1801.08112v5} (2023),
Theorem 3.16 and Corollary 3.21; the version indicated is the comparison text.
\bibitem{FJGAB} A. Gabrielov.
\emph{Multiplicity of a zero of an analytic function on a trajectory of a vector field}.
\href{https://arxiv.org/abs/math/9702229v1}{arXiv:math/9702229v1} (1997), Theorem 3.
\bibitem{FJBIN} G. Binyamini.
\emph{Multiplicity Estimates: a Morse-theoretic approach}.
Duke Mathematical Journal (2016).
\href{https://arxiv.org/abs/1406.1858v2}{arXiv:1406.1858v2} (2015), Corollary 1.
\bibitem{FJOPT} A. Ovchinnikov, G. Pogudin, and P. Thompson.
\emph{Parameter identifiability and input-output equations}.
\href{https://arxiv.org/abs/2007.14787v2}{arXiv:2007.14787v2} (2020),
Theorems 4.2 and 4.7 and Lemma 4.6.

\end{thebibliography}
\end{document}
